% CCSS 7.SP.A.2 — Making Inferences from Random Samples
\section{Making Inferences from Random Samples}

\begin{learningGoals}
\begin{itemize}
    \item Use data from random samples to make predictions about a population
    \item Understand variability in sample estimates
\end{itemize}
\end{learningGoals}

\begin{conceptBox}{Inference from Samples}
Predictions about a population are made using data from a random sample.
\end{conceptBox}

\begin{workedExample}{Estimating a Population Mean}
A random sample of $20$ students has an average test score of $82$. Predict the average for all $200$ students.
\answerHighlight{Estimate: $82$ is likely close to the population mean, but could vary.}
\end{workedExample}

\tipBox{Multiple random samples may give slightly different results.}

\begin{practiceBox}{Making Inferences}
\resetProblems

\prob A sample of $30$ plants has an average height of $15$ cm. Predict the average height for all $300$ plants.
\answerExplain{$15$ cm (estimate)}{The sample mean is a good estimate for the population mean.}

\prob Why do different samples give different estimates?
\answerExplain{Random variation}{Each sample may include different members, causing variation.}

\prob If two samples have means $20$ and $22$, what does this tell you?
\answerExplain{Population mean is likely between $20$ and $22$}{Multiple samples help narrow down the estimate.}

\prob How can you improve the reliability of an estimate?
\answerExplain{Use larger or more samples}{More data reduces the effect of random variation.}

\prob What is the risk of using a non-random sample?
\answerExplain{Biased prediction}{Non-random samples may not represent the population, leading to inaccurate predictions.}

\end{practiceBox}
